\section{Reconstrução do experimento}
\label{sec:corpus}

\subsection{Construção do corpus}

As URLs vieram do conjunto público \texttt{vinitvek/cybersecurityattacks}\footnote{\url{https://huggingface.co/datasets/vinitvek/cybersecurityattacks}}, filtrado por setor de saúde e pelo intervalo de 1 de janeiro a 30 de outubro de 2023. Os links serviram como índice para coletar, com Firecrawl, os textos completos nas páginas de origem, e não apenas os resumos do conjunto público. Das 329 URLs únicas selecionadas, 314 tiveram extração concluída. Após excluir 15 páginas inutilizáveis, restaram 299 notícias.

O original avaliou 60 notícias anotadas manualmente, limitadas aos primeiros 400 \textit{tokens}. Aqui, avaliamos textos completos com anotações automáticas revisadas manualmente. Portanto, mantivemos a tarefa, mas modificamos a seleção, a extensão dos textos e a produção das referências.

\subsection{Referências de anotação}

A anotação partiu de sugestões geradas pelo Muse Spark 1.3, um modelo de linguagem desenvolvido pela Meta \cite{meta2026muse}. Os anotadores revisaram manualmente os trechos e suas categorias, corrigiram marcações e acrescentaram entidades ausentes. Cada anotação delimita um trecho do texto e atribui uma categoria. Foram usadas quatro categorias. \texttt{IDTY} identifica organizações ou instituições. \texttt{ACT} corresponde a ações ou tipos de ataque. \texttt{VUL} marca vulnerabilidades. \texttt{LOC} identifica localidades geográficas.

A referência revisada, denominada \textit{overlay} nos artefatos, reúne 5.111 entidades nos 299 documentos. São 2.627 IDTY, 1.747 ACT, 65 VUL e 672 LOC. A referência de consenso contém 9.318 entidades e inclui menções repetidas da mesma entidade ao longo do texto. Na consolidação, foram mantidos os trechos da revisão final. Quando os mesmos limites apareciam nas outras anotações, os rótulos foram conciliados por voto majoritário. Em empate, prevaleceu o rótulo da revisão final. Trechos removidos ou corrigidos nessa revisão não foram reintroduzidos.

A cobertura não foi igual entre os anotadores. Dos trechos da referência de consenso, 4.640 têm três registros de anotação, 10 têm dois e 4.668 pertencem apenas à revisão final. Assim, as medidas de concordância entre os anotadores se referem aos trechos em comum, não à totalidade das 9.318 entidades.

Uma ferramenta local apoiou a conferência e a inclusão de entidades ausentes. Validadores automáticos verificaram rótulos e posições dos trechos, mas não garantem correção semântica. Nenhuma dessas referências equivale à anotação manual original.
